\bookchapter{Chapitre 58 : Aide-mémoire et ressources}

\chapterintro{Ce chapitre est un aide-mémoire utilisable de façon autonome. Quand tu ne te souviens plus d’une syntaxe ou que tu dois trancher vite, feuillete-le directement. Traite-le comme une « fiche de poche » : mémo de code courant, questions fréquentes, ressources d’apprentissage, pas une page de moins.}

\booksection{Les fiches mémo de code courant}

\begin{center}
\begin{tabularx}{\linewidth}{|X|X|}
\hline
\textbf{Ce que tu veux faire} & \textbf{Code clé} \tabularnewline
\hline
Mémoriser un état modifiable & var x by remember \{ mutableStateOf(valeurInitiale) \} \tabularnewline
Garder l’état à la rotation de l’écran & rememberSaveable \{ mutableStateOf(valeurInitiale) \} \tabularnewline
Disposition verticale/horizontale & \texttt{Column \{ \}}/\texttt{Row \{ \}} \tabularnewline
Liste à chargement paresseux & LazyColumn \{ items(liste) \{ ... \} \} \tabularnewline
Naviguer entre les écrans & navController.navigate("route") \tabularnewline
Gérer les tâches d’attente dans des coroutines & \texttt{viewModelScope.\allowbreak launch \{ .\allowbreak .\allowbreak .\allowbreak  \}} \tabularnewline
Observer un flux de données & \texttt{val list by flow.\allowbreak collectAsState()} \tabularnewline
Stocker une petite configuration & context.dataStore.edit \{ it[CLE] = valeur \} \tabularnewline
Stocker des données structurées & Le trio Room : @Entity / @Dao / @Database \tabularnewline
Récupérer des données en ligne & Interface Retrofit + fonctions suspend \tabularnewline
Envoyer une notification & NotificationChannel + NotificationCompat.Builder \tabularnewline
Écrire des logs & \texttt{Log.\allowbreak d("TAG", "message")} \tabularnewline
\hline
\end{tabularx}
\end{center}

\begin{center}\small\color{gray}Fiches mémo du développement Android (édition d’août 2026)\end{center}

\booksection{La foire aux questions (FAQ)}

Q1 : Je ne sais pas du tout programmer, combien de temps pour apprendre le développement Android ?

En suivant le parcours de ce livre : semaines 1-2, installer l’environnement et apprendre les bases de Kotlin ; semaines 3-4, apprendre Compose et créer une interface ; semaines 5-6, terminer le projet carnet de comptes ; semaines 7-8, finir le projet météo. Deux mois pour passer de zéro à une petite app faite tout seul. À condition d’y consacrer 1 à 2 heures par jour et de taper du code au lieu de seulement lire.

Q2 : Faut-il encore apprendre Java ?

Non. Nouveaux projets, nouveaux tutos, nouveaux postes : presque tout se fait en Kotlin. Tu ne rencontreras beaucoup de Java que pour « maintenir un projet vieux de dix ans ». Après Kotlin, si le travail l’exige, rattraper Java ira vite — les deux reposent sur le même écosystème de machine virtuelle.

Q3 : Émulateur ou appareil réel ?

L’émulateur est pratique pour développer et déboguer (snapshots, plusieurs modèles), mais la validation doit passer par un appareil réel — le tactile, les performances, le vrai réseau, les capteurs : l’émulateur ne les remplace pas. Si tu en as les moyens, acheter un appareil de milieu de gamme comme téléphone de test est l’investissement le plus rentable.

Q4 : Le développement Android, est-ce que ça trouve du travail ?

Le marché en 2026 : les postes juniors limités à « savoir faire des interfaces » sont très disputés, mais les développeurs « Kotlin solide + architecture + une app publiée » restent très recherchés. Pour toi qui apprends, ton meilleur CV, c’est cette app faite par toi et visible sur la boutique.

Q5 : Bloqué par un message d’erreur, que faire ?

Les trois étapes d’or : (1) recherche la première ligne de l’erreur telle quelle sur StackOverflow/dans un moteur de recherche ; (2) consulte la page correspondante de la documentation officielle ; (3) si ça ne marche toujours pas, envoie à l’IA ou à la communauté le paquet complet : « comportement attendu + comportement observé + message d’erreur intégral + code concerné ». La grande majorité des problèmes se règlent dès la première étape.

\booksection{Les ressources d’apprentissage recommandées}

\begin{center}
\begin{tabularx}{\linewidth}{|X|X|X|}
\hline
\textbf{Type} & \textbf{Ressource} & \textbf{Pourquoi c’est recommandé} \tabularnewline
\hline
Officiel & developer.android.com & La documentation officielle d’Android, tutos Compose/Kotlin compris : la référence numéro un \tabularnewline
Officiel & Le site officiel de Kotlin (kotlinlang.org) & Documentation officielle du langage + bac à sable en ligne \tabularnewline
Officiel & Google Codelabs & Des ateliers pas à pas très guidés, note maximale pour les débutants \tabularnewline
Communauté & Stack Overflow (tag android) & Le premier arrêt en cas d’erreur \tabularnewline
Communauté & Juejin / rubrique Android de CSDN & Beaucoup d’articles pratiques en chinois ; trie le neuf de l’ancien (regarde d’abord la date) \tabularnewline
Vidéo & La chaîne officielle Google Developers & Les explications des nouveautés présentées chaque année à la conférence des développeurs \tabularnewline
Compagnon & La septième partie « L’IA en pratique » (chapitres 37 à 47) & La méthode complète du développement Android assisté par l’IA \tabularnewline
\hline
\end{tabularx}
\end{center}

\begin{center}\small\color{gray}La liste des ressources d’apprentissage du développement Android (août 2026)\end{center}

\booksection{Aide-mémoire des commandes Gradle}

Au quotidien, les boutons suffisent, mais la ligne de commande est parfois plus efficace (surtout en la combinant avec des outils d’IA ou une chaîne CI) :

\begin{center}
\begin{tabularx}{\linewidth}{|X|X|}
\hline
\textbf{Commande} & \textbf{Rôle} \tabularnewline
\hline
\texttt{.\allowbreak /\allowbreak gradlew assembleDebug} & Construire un APK de débogage (sous Windows, retire le ./) \tabularnewline
\texttt{.\allowbreak /\allowbreak gradlew assembleRelease} & Construire la version release (signature configurée au préalable) \tabularnewline
\texttt{.\allowbreak /\allowbreak gradlew test} & Lancer tous les tests unitaires \tabularnewline
\texttt{.\allowbreak /\allowbreak gradlew lint} & Analyse statique : débusquer les bugs potentiels et le code peu orthodoxe \tabularnewline
\texttt{.\allowbreak /\allowbreak gradlew clean} & Nettoyer les artefacts de construction (la version compilée du « grand redémarrage ») \tabularnewline
\texttt{.\allowbreak /\allowbreak gradlew :app: dependencies} & Afficher l’arbre des dépendances (l’arme fatale contre les conflits de versions) \tabularnewline
\hline
\end{tabularx}
\end{center}

\begin{center}\small\color{gray}Aide-mémoire des commandes Gradle usuelles\end{center}

\begin{tcbinfo}{Astuce : en cas de souci de build, commence par clean}
Face aux phénomènes paranormaux du genre « hier ça tournait, aujourd’hui plus rien », six fois sur dix le cache de construction est en cause. Un ./gradlew clean suivi d’une reconstruction guérit la moitié de ces mystères. Si ça persiste, redémarre Android Studio — cette astuce se prend de préférence avec File Invalidate Caches.

\end{tcbinfo}

\booksection{Conclusion : de la première app à l’énième}

\booksection{Conclusion : de la première app à l’énième}

En regardant derrière toi : tu ne savais même pas ce qu’était une Activity, et te voilà avec l’environnement installé, Kotlin appris, des interfaces dessinées en Compose, une app reliée à sa base de données et au réseau, puis testée, empaquetée et publiée — sur ce chemin, tu as déjà franchi le plus dur du passage de zéro à un.

Le monde de la technique va très vite : les frameworks se mettent à niveau, les outils changent de génération, l’IA prend en charge de plus en plus de travail de « frappe clavier ». Mais une chose ne se périme pas : ta capacité à découper une idée floue en étapes claires, puis à la réaliser pas à pas. Les outils sont des amplificateurs, la réflexion est le socle. L’amplificateur se remplace, le socle ne bouge pas.

Tu n'es plus juste "quelqu'un qui sait se servir de son téléphone" : tu peux transformer les idées de ta tête en apps qui tournent réellement sur trois cents millions d'appareils Android. À partir d'aujourd'hui, crée avec ton clavier – c'est ça, la poésie du développeur Android.

\booksection{Références}

\begin{enumerate}
  \item Android Developers. (2026). Android Developer Guides and Jetpack Compose Documentation. developer.android.com.
  \item Kotlin Foundation. (2026). Kotlin Language Documentation. kotlinlang.org.
  \item Android Open Source Project. (2026). Android Platform Version History and API Levels.
  \item Google Play. (2026). Play Console Help: Publishing and App Signing.
  \item AndroidX Release Notes. (2026). Room / DataStore / Navigation / WorkManager releases.
  \item Square. (2026). Retrofit Documentation. github.com/square/retrofit.
\end{enumerate}
